\section{Primera etapa: adaptar los tres bloques de la robótica con los modelos fundacionales existentes}

La sección anterior planteó una ruta de dos etapas; esta sección examina la primera: incorporar los LLM, VLM y Code LM existentes al sistema robótico. Los tres bloques de la robótica tradicional son Planning, Perception y Actuation. Los modelos grandes sustituyen primero a la planificación, porque los modelos de lenguaje son buenos para descomponer una tarea en pasos; después, los VLM sustituyen a la percepción, porque pueden vincular la escena visual con la tarea expresada en lenguaje; y, un paso más allá, los Code LM intentan automatizar también la ejecución.

\subsection{SayCan: lo que puedo hacer, no lo que digo que haga}

El problema central de SayCan es este: el modelo de lenguaje propone muchas acciones, pero el robot no necesariamente puede realizarlas en el entorno actual. SayCan combina la descomposición de tareas del modelo de lenguaje con la affordance del robot, de modo que considera tanto «qué debería hacerse según el lenguaje» como «si el robot realmente puede hacerlo». Es un trabajo pionero en la conexión de los LLM con la planificación robótica.

\begin{figure}[H]
\centering
\includegraphics[width=0.96\textwidth]{ch07-01537.jpg}
\caption{SayCan: Query LLM to rank action primitives, y después se usa la affordance function para juzgar la viabilidad. Fotograma de la pantalla proyectada, aprox. 00:25:37.}
\end{figure}

\begin{knowledgebox}{Lectura de la figura: el LLM aporta el objetivo; la affordance, la restricción}
A la izquierda de la figura está la clasificación de acciones que hace el modelo de lenguaje; a la derecha, la evaluación de las posibilidades de acción del robot en el entorno real. Al leer esta figura conviene notar que SayCan no deja que el modelo de lenguaje controle directamente al robot, sino que añade una restricción entre el plan expresado en lenguaje y la viabilidad física.
\end{knowledgebox}

\begin{warningbox}{Un plan en lenguaje no equivale directamente a una acción del robot}
El robot debe considerar su capacidad de agarre, la posición en el espacio, el estado de los objetos y los fallos de ejecución. Si el LLM solo emite «toma la taza», no sabe si la taza está al alcance, si la pinza es adecuada ni si el entorno lo permite. El valor de SayCan está en aterrizar la intención expresada en lenguaje en la affordance.
\end{warningbox}

\subsection{Inner Monologue, DoReMi y VoxPoser}

Este apartado examina tres trabajos que siguen aplicando modelos fundacionales al razonamiento robótico. Inner Monologue subraya que, durante la ejecución, el robot necesita retroalimentación verbalizada y un monólogo interior que reincorpore la retroalimentación del entorno a la planificación. DoReMi se ocupa de cómo detectar y recuperarse cuando el plan y la ejecución no coinciden. VoxPoser, por su parte, usa el modelo de lenguaje para construir 3D value maps y convierte los objetivos expresados en lenguaje en campos de valor espaciales componibles, aplicados a la manipulation.

\begin{figure}[H]
\centering
\includegraphics[width=0.90\textwidth]{ch08-01649.jpg}
\caption{Inner Monologue: un ciclo de razonamiento corporizado formado por la planificación del modelo de lenguaje y la retroalimentación del entorno. Fotograma de la pantalla proyectada, aprox. 00:27:29.}
\end{figure}

\begin{knowledgebox}{Lectura de la figura: el robot necesita narrarse a sí mismo lo que ocurrió tras ejecutar}
La figura de Inner Monologue une Action, Environmental Feedback y Corrective Action en un ciclo cerrado. No se trata simplemente de que el robot «sepa hablar», sino de que el resultado de la ejecución regrese al proceso de planificación en lenguaje y forme un estado intermedio corregible.
\end{knowledgebox}

\begin{figure}[H]
\centering
\includegraphics[width=0.94\textwidth]{ch09-01737.jpg}
\caption{DoReMi: detectar y recuperarse de la plan-execution misalignment. Fotograma de la pantalla proyectada, aprox. 00:28:57.}
\end{figure}

\begin{knowledgebox}{Lectura de la figura: la detección de desajustes es un problema real de la inteligencia corporizada}
DoReMi trata el desajuste entre planificación y ejecución: el plan dice A, pero durante la ejecución el entorno se convierte en B, o el robot no completa la tarea como se esperaba. Nos recuerda que un sistema robótico no termina al generar un plan una sola vez, sino que debe detectar, corregir y recuperarse de forma continua.
\end{knowledgebox}

\begin{figure}[H]
\centering
\includegraphics[width=0.94\textwidth]{ch11-01956.jpg}
\caption{VoxPoser: usar el modelo de lenguaje para componer 3D value maps que guían la manipulación del robot. Fotograma de la pantalla proyectada, aprox. 00:32:36.}
\end{figure}

\begin{knowledgebox}{Lectura de la figura: del objetivo en lenguaje al campo de manipulación 3D}
La figura de VoxPoser convierte las instrucciones en lenguaje en un value map en el espacio. Al leerla hay que fijarse en la palabra «componible»: distintas restricciones pueden superponerse para formar un solo objetivo de manipulación, por ejemplo, acercarse al objeto, evitar obstáculos o mantener la postura.
\end{knowledgebox}

\subsection{De los Code LM a un verdadero modelo robótico}

La idea de que los Code LM sustituyan a Actuation consiste en que el modelo escriba código para el robot o llame a las API de bajo nivel. Sin embargo, Chen Jianyu (陈建宇) señala que esto todavía no es un robot foundation model completo. Se trata más bien de orquestar herramientas y controladores existentes; un verdadero VLA requiere que el propio modelo sea capaz de emitir acciones directamente a partir de la visión y el lenguaje, y no solo de escribir código de llamadas.

\begin{importantbox}{Los límites de la primera etapa}
Leveraging foundation models permite reforzar con rapidez la planificación, la percepción y la llamada a herramientas, pero en esencia sigue siendo un sistema modular. Un robot de propósito general necesita, en última instancia, entrenar con datos de robots un modelo capaz de manejar action.
\end{importantbox}

\subsection{Resumen de la sección}

La aportación de la primera etapa es demostrar que los LLM/VLM pueden ayudar al robot a comprender tareas, planificar acciones, detectar desajustes y construir objetivos espaciales. Sin embargo, estos métodos siguen dependiendo de los módulos robóticos tradicionales y todavía no constituyen un VLA de extremo a extremo.
